% Transcription — fonds Alexandre Grothendieck, shelfmark 117, batch 2 (pages 21-25).
% Established from the facsimile archives/batches/117/batch-02.pdf (a mirror of
% https://grothendieck.umontpellier.fr/117.pdf), per the skill
% transcribe-grothendieck. The batch file carries no cover sheet: PDF page N
% is page 20+N of the folder (the Montpellier footer prints « 21 » ... « 25 »),
% and the archivists' pencil numbers agree on every leaf (« 21 » ... « 25 »).
% The folder has 25 pages, so this batch is its last five.
%
% Pass: Opus 5.5 (claude-opus-5-5), 2026-09-23 — first pass, unchecked against the pages by a human.
%
% Three of the five pages are transcribed. Absent: pages 23 and 25, typed
% administrative documents (a D.E.A. notice of the USTL, a ministry circular
% with stamps) carrying nothing in his hand. Page 22 is written on the back of
% the leaf that became page 23 (the typescript shows through). He did not
% paginate these three leaves (batch 1 has his « 9 » on page 20, nothing here).
% Both \section headings are bracketed: they are ours, not his.
%
% What the batch is. Undated working notes on homotopy theories, felt pen.
%   21     Conditions on i : A_0 -> A so that ∫_{A_0} i^*(ξ) ⥲ ∫_A ξ for all
%          ξ ∈ H(A) (« cofinalité »); if i has a left adjoint j he presumes
%          i^*_H = j^H_! and deduces p_{0!} i^* ≃ p_!; example: A_0 a point
%          sent to a final object e, ∫_A ξ ≃ ξ(e); a footnote justifying the
%          circled formula by the square H(A) -> H(A^), H(B) -> H(B^) through
%          the Yoneda functors h_A, h_B; margin: « Cas des catégories
%          filtrantes ?? ». The second formula is written with i where the
%          argument needs j (noted in place, not corrected).
%   22     A Mayer-Vietoris / gluing condition: for two full subcategories
%          A', A'' of A, both sieves or both cosieves, union A, intersection
%          A_0, H(A) -> H(A') ×_{H(A_0)} H(A'') should be surjective on arrows
%          (« tendance d'être bijective sur les objets »), i.e. surjectivity
%          on Hom; to be used for the corner category (c -> a', c -> a''), its
%          dual, and products with objects B of Cat. Margin: A = A' ⊔_{A_0} A''
%          in (Cat).
%   24     A sheet of diagrams only: for f : X -> Y with Z -> X (g) and
%          h : ΩY -> Z, a first 3-step staircase, then a 7-row staircase of
%          squares with corners e running up to Ω^4 Y (labels Ω^k f, Ω^k g,
%          Ω^k h), the sequence Y <- X <- Z <- ΩY <- ΩX <- ... <- Ω^3 Y drawn
%          as a zigzag in which every composite of two consecutive arrows
%          factors through e, a perspective sketch of the staircase (described,
%          not redrawn), a dashed square, and a small P/Q sketch.
%
% Notation fixed for the batch: his script H as \mathcal{H}; the presheaf
% category as A^{\wedge}; his « Δ_0 » (bottom of the triangle, page 21) as
% \Delta_0; the base point e kept as e. Words he underlines are set in \emph.
\documentclass[11pt,a4paper]{article}
% Le préambule commun à toutes les transcriptions du fonds Grothendieck.
%
% Il tient en une idée : **l'incertitude se compose, elle ne se résout pas.**
% Une transcription qui lisse un mot douteux a détruit la seule chose pour
% laquelle on la faisait — un lecteur doit pouvoir savoir, à chaque endroit,
% ce qui était sur le feuillet et ce qui a été ajouté. Les six macros qui
% suivent sont donc l'appareil critique, et non de la décoration.
%
% Les mêmes commandes sont reconnues par `scripts/render.mjs`, qui produit la
% vue de lecture du site. Une macro ajoutée ici doit l'être aussi là-bas,
% faute de quoi le rendu échouera bruyamment — ce qui est voulu : un rendu
% silencieusement faux est pire qu'un rendu absent.

\ProvidesPackage{grothendieck}[2026/08/08 Transcriptions du fonds Grothendieck]

\usepackage{amsmath,amssymb,amsthm}
\usepackage{mathrsfs}
\usepackage{stmaryrd}
% Les diagrammes commutatifs sont partout dans ce fonds : c'est la langue
% même de la géométrie algébrique de Grothendieck, et une transcription qui
% les rendrait en prose ne transcrirait rien.
\usepackage{tikz-cd}
% Français par défaut — c'est la langue des feuillets — mais l'anglais est
% chargé aussi : la traduction et le résumé sont des documents anglais, et la
% typographie française (espaces avant les deux-points) y serait une faute.
% Ces documents basculent par \selectlanguage{english} en tête de corps.
\usepackage[english,french]{babel}
\usepackage{fontspec}
\usepackage{geometry}
\geometry{a4paper,margin=25mm,marginparwidth=28mm}
\usepackage{marginnote}
\usepackage{xcolor}
% ulem, not soul. `soul`'s \st and \ul are built on a character-by-character
% scan that XeLaTeX's font machinery breaks: the document fails with an
% undefined control sequence rather than degrading. ulem does the same two jobs
% and is Unicode-engine safe. `normalem` stops it hijacking \emph, which the
% transcriptions use for Grothendieck's own emphasis.
\usepackage[normalem]{ulem}
\usepackage{hyperref}
\hypersetup{colorlinks=true,linkcolor=black,urlcolor=black,pdfborder={0 0 0}}

% --- Métadonnées du lot -----------------------------------------------------
% Reprises telles quelles par le script de rendu, qui les affiche en tête de
% la vue de lecture. Elles fixent ce que le fichier prétend couvrir : sans
% elles, une transcription flotte sans point d'attache dans un fonds de seize
% mille feuillets.

\newcommand{\folder}[1]{\gdef\grothFolder{#1}}
\newcommand{\batch}[1]{\gdef\grothBatch{#1}}
\newcommand{\leaves}[2]{\gdef\grothFirst{#1}\gdef\grothLast{#2}}
\newcommand{\foldertitle}[1]{\gdef\grothFoldertitle{#1}}
\gdef\grothFolder{?}\gdef\grothBatch{?}\gdef\grothFirst{?}\gdef\grothLast{?}\gdef\grothFoldertitle{}

% --- Le résumé --------------------------------------------------------------
%
% Il ouvre la lecture modernisée, sous le titre « Résumé », et il est composé
% en petit. Ce n'est pas un ornement : c'est le seul passage du document qui
% ne soit pas des mathématiques, et un lecteur doit pouvoir le reconnaître
% avant de l'avoir lu — puis le sauter s'il connaît déjà le sujet, ou s'y
% arrêter s'il ne le connaît pas. Le filet qui le suit marque la reprise du
% corps.

\newenvironment{resume}
  {\small\setlength{\parskip}{0.45em}}
  {\par\medskip\hrule\medskip}

% --- Mots-clés ---------------------------------------------------------------
%
% En anglais, en clôture du résumé de la lecture modernisée : le vocabulaire
% moderne sous lequel chercher ce que les feuillets construisent. Le manifeste
% du site (`npm run manifest`) les extrait pour étiqueter la cote entière —
% cette ligne est la source unique des tags, il n'y a pas de fichier à côté.
\newcommand{\keywords}[1]{\par\smallskip\noindent
  {\footnotesize\textsc{Keywords} --- \textit{#1}}\par}

% --- L'appareil critique ----------------------------------------------------

% Le numéro de feuillet, en marge. C'est l'ancre qui permet de régler un
% désaccord sur une formule en regardant *un* feuillet plutôt qu'une chemise
% de six cents. Le site s'en sert aussi pour tourner les pages du fac-similé
% quand on fait défiler la transcription.
\newcommand{\leaf}[1]{%
  \marginnote{\footnotesize\color{gray!70}\textsf{#1}}%
  \label{leaf:#1}\ignorespaces}

% Illisible : on ne devine pas. Un mot inventé qui se lit comme les autres est
% le pire résultat possible.
\newcommand{\ill}{\textcolor{red!60!black}{[\,\ldots\,]}}

% Lecture douteuse : le mot est donné, et signalé comme incertain.
\newcommand{\uncertain}[1]{\uline{#1}}

% Ajout de l'éditeur — une lettre manquante, un mot sous-entendu.
\newcommand{\add}[1]{\textcolor{blue!50!black}{[#1]}}

% Biffé par Grothendieck lui-même. Ce qu'il a rayé fait partie du document :
% une rature dit souvent où la pensée a changé de direction.
\newcommand{\struck}[1]{\sout{#1}}

% Note du transcripteur, sur l'état matériel du feuillet.
\newcommand{\note}[1]{\par\smallskip{\small\color{gray!60!black}\textit{[#1]}}\par\smallskip}

% Note marginale de Grothendieck, distincte du corps du texte.
\newcommand{\marginal}[1]{\par\smallskip\noindent\hspace{1em}%
  {\small\color{gray!60!black}#1}\par\smallskip}

% --- Titre ------------------------------------------------------------------

\AtBeginDocument{%
  \begin{center}
    {\large\bfseries Fonds Alexandre Grothendieck --- cote n\textsuperscript{o}~\grothFolder}\\[2pt]
    {\itshape\grothFoldertitle}\\[4pt]
    {\small feuillets \grothFirst--\grothLast\ (lot \grothBatch)}\\[2pt]
    {\footnotesize\color{gray!60!black}Transcription établie d'après le fac-similé numérisé
     par l'Université de Montpellier.\\
     Les lectures incertaines sont soulignées, les illisibles notées [\,\ldots\,],
     les ajouts entre crochets.}
  \end{center}
  \vspace{1em}}

\endinput


\folder{117}
\batch{2}
\pages{21}{25}
\dating{1983}
\watermark{Édition de démonstration}
\foldertitle{Théories homotopiques : notes manuscrites (1983, s.d.).}

\begin{document}

\section{[Cofinalité, adjoints, et le recollement]}
\note{titre de l'éditeur, entre crochets : le feuillet 21 n'en porte aucun et
commence au milieu d'un développement (« Il faut aussi »). Les notations
$\mathcal{H}(A)$, $f^{*}$, $f^{\mathcal{H}}_{!}$, $\int_A$ sont celles des
feuillets précédents, au lot 1.}

\page{21}
Il faut aussi éventuellement des conditions sur un $i : A_0 \to A$ pour que
\[
  \int_{A_0} i^{*}(\xi) \xrightarrow{\ \sim\ } \int_A \xi
  \qquad \text{pour tout } \xi \in \mathcal{H}(A)
\]
\note{le $A$ but de $i$, et celui de $\mathcal{H}(A)$, sont repassés à l'encre
plus appuyée, peut-être sur une autre lettre ($B$ ?).}
(\uncertain{définition} \uncertain{de} « cofinalité »). Si $i$ a un adjoint
\uncertain{à} \struck{\uncertain{droite}} \struck{\ill{}} gauche $j$ (\ill{}
des foncteurs adjoints), je présume que
\[
  ?\qquad i^{*}_{\mathcal{H}} = j^{\mathcal{H}}_{!}\ {}^{(*)},
\]
\note{la formule est cernée d'un ovale, précédée d'un point d'interrogation ;
l'astérisque renvoie à la note « * Résulte… » du bas du feuillet.}
et la formule sort, i.e.
\[
  p^{\mathcal{H}}_{0!}\, i^{*}_{\mathcal{H}} \xrightarrow{\ \sim\ }
  p^{\mathcal{H}}_{!}
\]
\uncertain{dérivant}
\[
  p^{\mathcal{H}}_{0!}\, i^{\mathcal{H}}_{!} \simeq p^{\mathcal{H}}_{!}
\]
qui est conséquence de $p_0 i = p$.
\note{sur la page, les deux dernières formules portent $i$ ; or
$p_0 : A_0 \to \Delta_0$ et $i : A_0 \to A$ ne se composent pas en $p_0 i$,
ni $p^{\mathcal{H}}_{0!}$ avec $i^{\mathcal{H}}_{!}$. L'argument demande
l'adjoint $j : A \to A_0$, avec $j^{\mathcal{H}}_{!} = i^{*}_{\mathcal{H}}$ et
$p_0 j = p$. Le signe est peut-être un $j$ tracé sans sa queue ; on transcrit
ce qu'on lit.}

À gauche des formules, le triangle :

\begin{tikzcd}
A_0 \arrow[r, "i"] \arrow[d, "p_0"'] & A \arrow[dl, no head, "p"] \\
\Delta_0 &
\end{tikzcd}

\note{l'oblique $p$ est un trait sans pointe. Le but, noté $\Delta_0$, est
écrit à double trait, comme un delta ; on le lit comme la catégorie ponctuelle,
ce que la page ne dit pas.}

\emph{Ex} Cas où $A_0$ réduit à un pt, envoyé dans l'objet final $e$ de $A$ :
on trouve
\[
  \int_A \xi \xleftarrow{\ \sim\ } \xi(e)
\]
\note{la flèche est repassée plusieurs fois ; on la lit comme une flèche vers
la gauche marquée $\sim$, sans certitude sur son sens.}

\marginal{Cas des catégories filtrantes ?? (« \uncertain{Ets} » \ill{}
\uncertain{cofinal} \ill{} \uncertain{valeurs})}
\note{note écrite en oblique dans la marge gauche, en face des premières
lignes ; la fin, sous parenthèse, est à peine lisible.}

${}^{(*)}$ Résulte \uncertain{du} diagramme commutatif (\ill{} en remplaçant
$A_0$ par $B$), valable pour $j : A \to B$ quelc. :

\begin{tikzcd}
\mathcal{H}(A) \arrow[r, "(h_A)^{\mathcal{H}}_{!}"] \arrow[d, "j^{\mathcal{H}}_{!}"'] & \mathcal{H}(A^{\wedge}) \arrow[d, "(j^{*}_{\wedge})^{*}_{\mathcal{H}}"] \\
\mathcal{H}(B) \arrow[r, "(h_B)^{\mathcal{H}}_{!}"'] & \mathcal{H}(B^{\wedge})
\end{tikzcd}

\emph{NB} dans le cas où $j$ a un \struck{\ill{}} adjoint à dr. $i$,
\uncertain{on a} $j^{*}_{\wedge} = i^{\wedge}_{!}$, qui induit
$i : B \to A$, donc $(i^{\wedge}_{!})^{*}_{\mathcal{H}}$ induit
$i^{*}_{\mathcal{H}}$.

\marginal{diagrammes \ill{} \struck{\ill{}} : \ill{}}
\note{note encadrée au bas de la marge gauche, reliée par un trait au
diagramme ; au-dessous, un « \uncertain{II} » souligné, suivi d'un trait
horizontal, sans rien après.}

\page{22}
On veut que pour deux sous-catégories pleines $A'$, $A''$ de $A$
\uncertain{de réunion} $A$, qui \uncertain{soient} soit des cribles, soit des
cocribles, de réunion $A$, d'intersection $A_0$, \uncertain{d'injection}, \struck{$\mathrm{Hom}($}
\[
  \mathcal{H}(A) \longrightarrow
  \mathcal{H}(A') \times_{\mathcal{H}(A_0)} \mathcal{H}(A'')
\]
soit \uncertain{donc} \emph{surjective} sur les flèches (elle a tendance
d'être \uncertain{bijective} sur les objets …), i.e. pour
$\xi, \eta \in \mathrm{Ob}\,\mathcal{H}(A)$, on veut
\[
  \mathrm{Hom}(\xi, \eta) \longrightarrow
  \mathrm{Hom}(\xi', \eta') \times_{\mathrm{Hom}(\xi_0, \eta_0)}
  \mathrm{Hom}(\xi'', \eta'')
\]
\emph{surjectif}.
\note{« de réunion $A$ » paraît deux fois, à la deuxième et à la quatrième
ligne ; la première occurrence est une lecture douteuse. « surjectif » est
souligné deux fois.}

\marginal{\emph{NB} On aura alors $A = A' \amalg_{A_0} A''$ dans (Cat)}

\marginal{(\uncertain{c'est} \uncertain{$J$})}
\note{en marge gauche, en face de la seconde formule ; lecture très douteuse.}

On l'utilisera dans le cas de \struck{$A$ \ill{}}

\begin{tikzcd}
\bullet & \\
c \arrow[u, "A'"] \arrow[r, "A''"'] & \bullet
\end{tikzcd}

\note{les deux flèches issues de $c$ sont cernées chacune d'un ovale, marqué
$A'$ (la verticale) et $A''$ (l'horizontale) ; les deux autres sommets ne
sont pas nommés.}

et leur cas duals, et le produit par \uncertain{les} éléments $B$ de Cat

\section{[Diagrammes : la suite des lacets]}
\note{titre de l'éditeur, entre crochets : le feuillet 24 ne porte que des
diagrammes, sans un mot.}

\page{24}
\begin{tikzcd}[column sep=small, row sep=small]
 & e \arrow[d] & \Omega Z \arrow[l] \arrow[d] \\
e \arrow[d] & \Omega Y \arrow[l] \arrow[d] & \Omega X \arrow[l] \arrow[d] \\
X \arrow[d] & Z \arrow[l] \arrow[d] & e \arrow[l] \\
Y & e \arrow[l] &
\end{tikzcd}

\begin{tikzcd}[column sep=small, row sep=small, nodes={font=\scriptsize}]
 & & & & e \arrow[d] & \Omega^{3} Z \arrow[l] \arrow[d, "\Omega^{3} g"] & \Omega^{4} Y \arrow[l, "\Omega^{3} h"'] \arrow[d] \\
 & & & e \arrow[d] & \Omega^{3} Y \arrow[l] \arrow[d, "\Omega^{2} h"] & \Omega^{3} X \arrow[l, "\Omega^{3} f"'] \arrow[d] & e \arrow[l] \\
 & & e \arrow[d] & \Omega^{2} X \arrow[l] \arrow[d, "\Omega^{2} f"] & \Omega^{2} Z \arrow[l, "\Omega^{2} g"'] \arrow[d] & e \arrow[l] & \\
 & e \arrow[d] & \Omega Z \arrow[l] \arrow[d, "\Omega g"] & \Omega^{2} Y \arrow[l, "\Omega h"'] \arrow[d] & e \arrow[l] & & \\
e \arrow[d] & \Omega Y \arrow[l] \arrow[d, "h"] & \Omega X \arrow[l, "\Omega f"'] \arrow[d] & e \arrow[l] & & & \\
X \arrow[d, "f"'] & Z \arrow[l, "g"] \arrow[d] & e \arrow[l] & & & & \\
Y & e \arrow[l] & & & & &
\end{tikzcd}

\note{les exposants de $\Omega^{2} g$ et $\Omega^{2} f$ sont un petit trait
qu'on lit « 2 » ; celui de $\Omega^{3} f$ est tracé comme les « 3 » voisins.
Le sommet est écrit $\Omega^{4} Y$ (un « 4 » peu formé).}

À droite, isolé :

\begin{tikzcd}
 & P \arrow[d, no head] \\
P & Q \arrow[l]
\end{tikzcd}

\note{lettres douteuses : le $P$ du haut est repassé ; le trait vertical n'a
pas de pointe nette.}

\begin{tikzcd}[column sep=small, row sep=small, nodes={font=\scriptsize}]
 & & e \arrow[dl] & & e \arrow[dl] & & e \arrow[dl] & & e \arrow[dl] & & \\
Y & X \arrow[l] & Z \arrow[l] \arrow[dl] & \Omega Y \arrow[l] \arrow[ul] & \Omega X \arrow[l] \arrow[dl] & \Omega Z \arrow[l] \arrow[ul] & \Omega^{2} Y \arrow[l] \arrow[dl] & \Omega^{2} X \arrow[l] \arrow[ul] & \Omega^{2} Z \arrow[l] \arrow[dl] & \Omega^{3} Y \arrow[l] \arrow[ul] & \cdots \arrow[l] \\
 & e \arrow[ul] & & e \arrow[ul] & & e \arrow[ul] & & e \arrow[ul] & & &
\end{tikzcd}

\note{la ligne du milieu, $Y \leftarrow X \leftarrow Z \leftarrow \Omega Y
\leftarrow \cdots \leftarrow \Omega^{3} Y \leftarrow$, est écrite une première
fois seule, plus haut, jusqu'à $\Omega Y \leftarrow$, suivie de deux points ;
elle est reprise ici en zigzag, chaque composé de deux flèches consécutives
passant par un $e$, alternativement au-dessus et au-dessous. La flèche
$e \to Y$ de gauche est surchargée.}

\begin{tikzcd}[column sep=small, row sep=small]
e \arrow[d] & \cdot \arrow[l, dashed] \arrow[dd, dashed, no head] & \bullet \arrow[l, dashed, no head] \arrow[d, dashed, no head] \\
X \arrow[d] & & e \arrow[ll] \arrow[dl] \\
Y & e \arrow[l] &
\end{tikzcd}

\note{les traits pointillés sont ceux de la page : la flèche du haut vers $e$,
le trait vertical qui descend du milieu de celle-ci vers le $e$ du bas, et le
trait en pointillé serré qui relie le coin supérieur droit au $e$ de droite
(son extrémité basse est peut-être une pointe).}

\note{à gauche de ce carré, un croquis en perspective de l'escalier
précédent : un réseau de flèches dont les sommets $e$ sont cerclés, $X$ et $Y$
en bas à gauche, traversé par un trait épais en escalier et par un axe
vertical ; on ne le redessine pas.}

\end{document}
